% Coe College summer REU report template.
%
% Edit this file. Leave coereu.cls alone unless you're changing the design.
% See README.md for the full command reference.
%
% Build: latexmk -pdf main.tex
%   or:  pdflatex -> biber -> pdflatex -> pdflatex
% Compile with pdfLaTeX, NOT XeLaTeX.
%
% Based on "A soft template for homework solutions" by Lucas R. Ximenes
% (Jimeens). Retained with attribution as that template's license asks.
%%%%%%%%%%%%%%%%%%%%%%%%%%%%%%%%%%%%%%%%%%%%%%%%%%%%%%%%%%%%%%%%%

\documentclass{coereu} % letterpaper, twoside, 11pt -- pass options here, e.g. [oneside]

% Delete this line once bib.bib holds real references and you want a
% References section. Until then it keeps an empty one from appearing.
\coeprintbibfalse

\begin{document}

%%%%%%%%%%%%%%%%%%%%%%%%%%%%%%%%%%%%%%%%%%%%%%%%%%%%%%%%%%%%%%%%%
% COVER -- title, subtitle, your name
%%%%%%%%%%%%%%%%%%%%%%%%%%%%%%%%%%%%%%%%%%%%%%%%%%%%%%%%%%%%%%%%%
\pretitle
{Your Project Title Here}          % ⟸ Main title
{A Coe College Summer Research Report} % ⟸ Subtitle
{Your Name}                        % ⟸ Your name

% Report metadata. All have defaults -- override only what you need.
\renewcommand{\reportNumber}{1}                        % big chapter numeral
\renewcommand{\reportKicker}{Coe College REU}          % spine label + header
\renewcommand{\coeProgram}{Research at Coe College}    % cover + footer line

%%%%%%%%%%%%%%%%%%%%%%%%%%%%%%%%%%%%%%%%%%%%%%%%%%%%%%%%%%%%%%%%%
% BODY
% \startcontents must stay before your first \chapter -- it feeds the
% mini-summary box on the chapter page.
%%%%%%%%%%%%%%%%%%%%%%%%%%%%%%%%%%%%%%%%%%%%%%%%%%%%%%%%%%%%%%%%%
\makeatletter
    \startcontents[sections]
    \phantomsection
    \chapter{Summer Research Report}
\makeatother

\section{Introduction}\label{sec:intro}

State what problem you worked on and why it matters, in language a
scientifically literate reader outside your subfield can follow. Say what
the report covers, and close with a sentence or two mapping out the
sections that follow.

Use \verb|\divider| to separate passages within a section:

\divider

Refer to other sections by name with \verb|\nameref|, like
\nameref{sec:methods} --- not \verb|\ref|, which returns the wrong thing
because sections here are deliberately unnumbered. Figures and equations
\emph{are} numbered, so \verb|\ref| is correct for those.

\section{Background}\label{sec:background}

The context a reader needs before your own work makes sense: the state of
the field, the instrument or system you worked with, and the open question
your project addresses.

Cite sources with \verb|\cite{key}| after adding entries to
\verb|bib.bib|, or with \verb|\textcite{key}| when the author is part of
your sentence.

\section{Methods}\label{sec:methods}

What you actually did. Be concrete about tools, techniques, and choices ---
including the ones that did not work. A short, honest account of what broke
and how you diagnosed it is usually more valuable to the reader than a
tidy narrative in which everything worked the first time.

To include a figure:

\begin{verbatim}
\begin{figure}[H]
    \centering
    \includegraphics[width=0.8\textwidth]{yourfigure.png}
    \caption{What the reader should notice.}
    \label{fig:example}
\end{figure}
\end{verbatim}

\section{Results}\label{sec:results}

What you found. Use the \verb|results| environment to set off a finding you
want the reader to stop at:

\begin{results}[optional label]
    A highlighted passage. The box breaks across pages on its own, so it
    can hold as much as you need.
\end{results}

For a single equation that matters more than the rest, use
\verb|keyresult|, which boxes it:

\begin{keyresult}\label{eq:example}
    E = mc^{2}
\end{keyresult}

\section{Discussion}\label{sec:discussion}

What the results mean, how far they can be trusted, and what limits them.
Say plainly where the uncertainty lies --- a stated limitation reads as
competence, not weakness.

\section{Conclusions and Future Work}\label{sec:conclusions}

Summarise briefly, then say what comes next: the questions still open and
what you would do with more time.

\section{Acknowledgements}\label{sec:ack}

Thank your advisor, collaborators, and anyone whose work you built on.
Include the funding acknowledgement here --- confirm the exact grant
number with your advisor rather than copying it from another document.

\end{document}
